\section{Related Work}
\label{sec:related_work}

Our proposed Zero-Shot Calibration-Free Model Merging framework is situated at the intersection of parameter merging, dynamic activation-space ensembling, and test-time adaptation (TTA). In this section, we contextualize our calibration-free paradigms (EER and EPL-OCA) relative to recent literature, highlighting their practical and operational advantages.

\subsection{LoRA and Weight-Space Parameter Merging}
Low-Rank Adaptation (LoRA) \cite{hu2021lora} has become the dominant paradigm for parameter-efficient fine-tuning (PEFT), allowing task-specific low-rank updates to be applied to frozen pre-trained backbone models. While LoRA reduces training costs, serving dozens of distinct LoRA adapters in production can introduce severe memory bottlenecks if experts must be physically swapped in and out of GPU/CPU memory on-the-fly. To circumvent this, researchers have proposed static weight-space merging, including Task Arithmetic \cite{ilharco2022editing}, TIES-Merging \cite{yadav2023ties}, and AdaMerging \cite{adamerging}. These methods compress multiple adapters into a single static set of weights. 

However, static weight merging suffers from three fundamental limitations from a deployment perspective: (1) \emph{Parameter Interference:} merging orthogonal or conflicting tasks (such as natural photos vs. handwritten digits) in weight-space frequently corrupts specialized representations, leading to severe performance degradation. (2) \emph{Homogeneous Limitation:} they assume all inputs are uniformly distributed or require the merged model to support all tasks simultaneously, compromising specialized expert performance. (3) \emph{Iterative Optimization Cost:} methods such as AdaMerging require expensive offline optimization loops to determine task-merging weights, which cannot adapt to streaming, dynamic real-world shifts.

\subsection{Dynamic Routing and Activation-Space Ensembling}
To overcome the limitations of static merging, dynamic ensembling methods execute expert adapters concurrently in the activation space or route input tokens dynamically. Mixture-of-Experts (MoE) \cite{shazeer2017outrageously} utilizes learned routing gates, but requires expensive joint training from scratch, which is impractical when combining pre-existing independent experts. 

Recently, training-free, single-pass activation ensembling frameworks like SABLE \cite{sable} and SPS-ZCA \cite{sps_zca} have been introduced. These methods route inputs by projecting test-time representations onto task-specific centroid anchors pre-computed from a small calibration split. PFSR \cite{pfsr} uses frozen expert classification heads as task anchors. While these methods are training-free and fast, they all depend on \textbf{offline, labeled calibration datasets} ($|\mathcal{C}_k|=64$) to locate task centroids before deployment. Under tight data privacy constraints or out-of-the-box serving where task labels are completely missing, these calibration-dependent methods are inapplicable. Our work, in contrast, is the first to achieve dynamic activation ensembling completely unsupervised and calibration-free, directly addressing this real-world deployment barrier.

\subsection{Test-Time Adaptation (TTA)}
Test-time adaptation (TTA) aims to adjust a pre-trained model to non-stationary test distributions or domain shift during deployment. Classic methods such as TENT \cite{wang2020tent} and CoTTA \cite{wang2022continual} update model parameters (such as batch normalization parameters or adapter weights) by minimizing prediction entropy on incoming test batches. 

While TTA can handle covariate shifts, it is heavily bottlenecked on edge hardware due to its dependence on \textbf{online backpropagation}. Computing backward-pass gradients during inference:
\begin{enumerate}
    \item Generates massive memory overhead (requiring the device to store activation graphs in RAM), which can easily trigger Out-Of-Memory (OOM) crashes on low-power microcontrollers.
    \item Significantly inflates serving latency (often by $3\times$ to $5\times$), defeating the strict real-time constraints of robotics and streaming applications.
    \item Suffers from optimization instability, where continuous gradient updates on corrupted inputs can cause representation collapse or catastrophic forgetting.
\end{enumerate}
Our proposed calibration-free ensembling paradigms operate completely in the forward pass. They use training-free, geometrically-sound entropy pseudo-labeling and exponential moving average (EMA) centroid tracking to adapt to covariate shift. They require \textbf{zero backward passes}, maintaining low, deterministic latency and stable on-device serving.
